\begin{abstract}

本文主要研究了微网购电与储能问题。光伏出力与负载需求的偏差会增加微网补购电量，储能与合同调整则为降低费用提供了空间。本文以购电费用最小为目标，建立含储能的线性规划模型，依次研究确定性调度、紧急购电、合同调整和波动电价。\par\medskip

针对问题一，我们将负载与光伏预测换算为10分钟时段电量，以能量守恒、储电量递推与边界、充放电功率上限及首末储电量相等为约束，在此基础上建立了确定性线性规划模型，并采用HiGHS求解。
求解结果如下：

\textbf{全天购电量为59482.70~kWh，购电费为35126.95元，较无储能方案节省26.90\%。}\par\medskip

针对问题二，我们根据周同期负载和近期光伏数据预测净负荷并校正偏差，按误差分位数增加采购保护量，结合末端储电价值确定日前合同。
在此基础上建立了含紧急购电惩罚的滚动计划模型。
求解结果如下：

执行时根据实际余缺充放电，储能不足的部分由紧急购电补足。\textbf{统计期内购电费为13790848.69元，其中紧急购电费为468467.64元，占3.40\%。}\par\medskip

针对问题三，我们将当时已发布的光伏预报细化为10分钟曲线，与历史统计预测按误差加权组合；由线性规划生成备选合同，再按历史误差情景中的净增益均值决定是否调整。
在此基础上建立了基于线性规划与情景净增益的滚动优化模型。
求解结果如下：

\textbf{统计期内购电费为13298484.49元，较问题二方案节省3.57\%，紧急购电量下降43.43\%，1002次机会中调整265次。}日前组合光伏预测的MAE为127.90~kW，较历史预测降低14.83\%。\par\medskip

针对问题四，我们用前一周同时段电价与历史偏差中位数预测价格，分别求解合同不可调整与可调整两类场景。可调整场景按同一历史日期配对净负荷误差和对数价格误差，比较备选合同的情景净增益均值。
在此基础上建立了价格误差情景驱动的购电储能优化模型。
求解结果如下：

\textbf{两类方案购电费分别为14526595.92元和14018361.30元，较固定电价分别上升5.34\%和5.41\%；可调整方案较不可调整方案节省3.50\%。}\par\medskip

本文的创新主要体现在预测信息利用与合同调整评价的结合。针对单一预测难以及时反映光伏变化、合同调整难以兼顾当期费用与跨日储电价值的问题，我们采用分时发布预报与历史预测的误差加权组合，依据历史净负荷误差分时设置采购保护量，并在相同决策点、储电状态和历史误差路径下配对比较调整与保持合同，将调整费用及末端储电短缺机会成本纳入净增益评价。\textbf{该方法提升了预测信息的利用效率，增强了合同修正对供需变化的适应性，并使费用评价兼顾当前结算与跨日储能价值，可用于微网日前购电、日内合同修正及虚拟电厂储能协同调度。}

\keywords{微网经济调度；线性规划；滚动时域优化；组合预测；经验分位数；情景分析}

\end{abstract}
